% ----------------------------------------------------------
% Apendices - Configuracoes YAML (Blindado contra erros)
% ----------------------------------------------------------

\begin{apendicesenv}

% ---
\chapter{Dados brutos da caracterização do sensor de corrente}
\label{apendice:ensaio_hall}
% ---

Coleta completa do ensaio descrito na \autoref{sec:same_sensor_corrente}, cujos
resultados são discutidos na \autoref{sec:resultado_sensor_corrente}. Medem-se
$I$ e $V_{out}$; $B$ vem da \autoref{eq:B_hopkinson} e
$\Delta V = V_{out} - V_O$ usa o intercepto da regressão de cada configuração.

\begin{table}[H]
    \centering
    \caption{Coleta completa da caracterização estática do sensor de corrente}
    \label{tab:hall_dados_brutos}
    \small
    \begin{tabular}{cccccc}
        \toprule
        \textbf{$N$} & \textbf{$I$ [\si{\ampere}]} &
        \textbf{$V_{out}$ [\si{\volt}]} & \textbf{$\Delta V$ [\si{\milli\volt}]} &
        \textbf{$B$ [G]} & \textbf{Resíduo [\si{\milli\volt}]} \\
        \midrule
        1 & 0{,}000 & 2{,}48570 & $-$0{,}945 & 0{,}00 & $-$0{,}945 \\
        1 & 0{,}159 & 2{,}48975 & 3{,}105 & 2{,}47 & 0{,}124 \\
        1 & 1{,}040 & 2{,}50530 & 18{,}655 & 16{,}12 & $-$0{,}844 \\
        1 & 1{,}539 & 2{,}51590 & 29{,}255 & 23{,}86 & 0{,}400 \\
        1 & 2{,}012 & 2{,}52550 & 38{,}855 & 31{,}19 & 1{,}131 \\
        1 & 2{,}529 & 2{,}53520 & 48{,}555 & 39{,}21 & 1{,}138 \\
        1 & 2{,}939 & 2{,}54200 & 55{,}355 & 45{,}57 & 0{,}250 \\
        1 & 3{,}113 & 2{,}54490 & 58{,}255 & 48{,}26 & $-$0{,}112 \\
        1 & 4{,}082 & 2{,}56270 & 76{,}055 & 63{,}29 & $-$0{,}481 \\
        1 & 5{,}041 & 2{,}58050 & 93{,}855 & 78{,}16 & $-$0{,}661 \\
        \midrule
        4 & 0{,}000 & 2{,}48640 & $-$1{,}232 & 0{,}00 & $-$1{,}232 \\
        4 & 0{,}101 & 2{,}49780 & 10{,}168 & 6{,}26 & 0{,}990 \\
        4 & 0{,}355 & 2{,}52050 & 32{,}868 & 22{,}02 & 0{,}610 \\
        4 & 0{,}620 & 2{,}54420 & 56{,}568 & 38{,}45 & 0{,}231 \\
        4 & 1{,}056 & 2{,}58320 & 95{,}568 & 65{,}49 & $-$0{,}387 \\
        4 & 1{,}662 & 2{,}63800 & 150{,}368 & 103{,}07 & $-$0{,}652 \\
        4 & 2{,}094 & 2{,}67800 & 190{,}368 & 129{,}86 & 0{,}094 \\
        4 & 2{,}522 & 2{,}71710 & 229{,}468 & 156{,}40 & 0{,}303 \\
        4 & 2{,}766 & 2{,}73900 & 251{,}368 & 171{,}53 & 0{,}032 \\
        4 & 3{,}034 & 2{,}76340 & 275{,}768 & 188{,}16 & 0{,}080 \\
        4 & 4{,}038 & 2{,}85450 & 366{,}868 & 250{,}42 & $-$0{,}050 \\
        4 & 5{,}077 & 2{,}94880 & 461{,}168 & 314{,}85 & $-$0{,}160 \\
        4 & 5{,}329 & 2{,}97200 & 484{,}368 & 330{,}48 & 0{,}142 \\
        \midrule
        8 & 0{,}000 & 2{,}49150 & 0{,}801 & 0{,}00 & 0{,}801 \\
        8 & 0{,}580 & 2{,}58630 & 95{,}601 & 71{,}94 & $-$0{,}341 \\
        8 & 1{,}022 & 2{,}65960 & 168{,}901 & 126{,}76 & $-$0{,}155 \\
        8 & 1{,}514 & 2{,}74090 & 250{,}201 & 187{,}78 & $-$0{,}241 \\
        8 & 2{,}014 & 2{,}82410 & 333{,}401 & 249{,}80 & 0{,}251 \\
        8 & 2{,}989 & 2{,}98440 & 493{,}701 & 370{,}73 & $-$0{,}731 \\
        8 & 3{,}319 & 3{,}03893 & 548{,}231 & 411{,}66 & $-$0{,}788 \\
        8 & 4{,}164 & 3{,}18070 & 690{,}001 & 516{,}47 & 1{,}204 \\
        \bottomrule
    \end{tabular}

    \vspace{2pt}
    {\small Fonte: Autor (2026)}
\end{table}

% ---
\chapter{Dados brutos da calibração do tubo de Pitot}
\label{apendice:ensaio_pitot}
% ---

Coleta completa do ensaio em túnel de vento descrito na
\autoref{sec:calibracao_pitot}. O tubo de Pitot instrumentado com o sensor de
pressão diferencial MPXV7002DP foi lido por um conversor A/D de 10 bits com
referência de \SI{5}{\volt}, em seis patamares de velocidade entre
\SI{1}{\meter\per\second} e \SI{14,8}{\meter\per\second}, com aquisição
contínua de 55 a 186 amostras por patamar.

\begin{table}[H]
    \centering
    \caption{Condições do ensaio de calibração do tubo de Pitot}
    \label{tab:pitot_condicoes}
    \small
    \begin{tabular}{lcc}
        \toprule
        \textbf{Parâmetro} & \textbf{Valor} & \textbf{Unidade} \\
        \midrule
        Densidade do ar ($\rho$)   & 1,110 & \si{\kilo\gram\per\meter\cubed} \\
        Temperatura                & 25    & \si{\celsius} \\
        Pressão atmosférica        & 95,0  & \si{\kilo\pascal} \\
        Resolução do conversor     & 10    & bits \\
        Tensão de referência       & 5,0   & \si{\volt} \\
        Resolução equivalente      & 4,89  & \si{\milli\volt}/LSB \\
        \bottomrule
    \end{tabular}

    \vspace{2pt}
    {\small Densidade calculada por $\rho = P/(R\,T)$ a partir da pressão e da
    temperatura acima; não foi medida durante o ensaio. Fonte: Autor (2026)}
\end{table}

\begin{table}[H]
    \centering
    \caption{Coleta por patamar de velocidade}
    \label{tab:pitot_dados_brutos}
    \small
    \begin{tabular}{cccccc}
        \toprule
        \textbf{$v_{ref}$} & \textbf{Amostras} & \textbf{ADC médio} &
        \textbf{Desvio} & \textbf{Amplitude} & \textbf{Tensão média} \\
        {[\si{\meter\per\second}]} & & [LSB] & [LSB] & [LSB] & [\si{\volt}] \\
        \midrule
        1,00  & 55  & 546,02 & 0,216 & 1,2  & 2,6687 \\
        2,89  & 155 & 545,38 & 1,959 & 16,6 & 2,6657 \\
        5,60  & 159 & 547,20 & 0,886 & 6,9  & 2,6745 \\
        8,40  & 149 & 549,86 & 0,775 & 7,1  & 2,6875 \\
        11,57 & 155 & 554,51 & 0,868 & 4,7  & 2,7103 \\
        14,80 & 186 & 561,07 & 2,445 & 24,1 & 2,7424 \\
        \bottomrule
    \end{tabular}

    \vspace{2pt}
    {\small Fonte: Autor (2026)}
\end{table}

A excursão total do conversor em toda a faixa ensaiada é de apenas
\SI{15,7}{LSB}, ou cerca de \SI{77}{\milli\volt}: o sensor trabalha numa
janela estreita da escala. O ajuste por mínimos quadrados do modelo
quadrático $ADC = k\,v^2 + ADC_0$ fornece
$k = \num{0,071521}$~LSB/(\si{\meter\per\second})$^2$ e
$ADC_0 = \num{545,14}$~LSB, com $R^2 = \num{0,9944}$; a
\autoref{tab:pitot_verificacao} confronta a velocidade reconstruída por
$v = \sqrt{(ADC - ADC_0)/k}$ contra a referência.

\begin{table}[H]
    \centering
    \caption{Verificação do ajuste quadrático}
    \label{tab:pitot_verificacao}
    \small
    \begin{tabular}{ccccc}
        \toprule
        \textbf{$v_{ref}$ [\si{\meter\per\second}]} & \textbf{ADC [LSB]} &
        \textbf{$v$ calculada [\si{\meter\per\second}]} &
        \textbf{Erro [\si{\meter\per\second}]} & \textbf{Erro [\si{\percent}]} \\
        \midrule
        1,00  & 546,02 & 3,51  & $+$2,51 & 250,8 \\
        2,89  & 545,38 & 1,83  & $-$1,06 & 36,6 \\
        5,60  & 547,20 & 5,37  & $-$0,23 & 4,2 \\
        8,40  & 549,86 & 8,12  & $-$0,28 & 3,3 \\
        11,57 & 554,51 & 11,45 & $-$0,12 & 1,1 \\
        14,80 & 561,07 & 14,92 & $+$0,12 & 0,8 \\
        \bottomrule
    \end{tabular}

    \vspace{2pt}
    {\small Fonte: Autor (2026)}
\end{table}

Acima de \SI{5}{\meter\per\second} o erro médio fica em \SI{2,3}{\percent} e
cai para menos de \SI{1}{\percent} no topo da faixa. Nos dois primeiros
patamares o erro é inaceitável, e a causa não é o ajuste: entre 0 e
\SI{3}{\meter\per\second} a leitura varia menos de \SI{1}{LSB}, de modo que a
quantização e o ruído dominam completamente o sinal.

As constantes $k$ e $ADC_0$ acima estão amarradas simultaneamente ao sensor,
ao ar e ao conversor, e por isso não são transportáveis. Separando os três
estágios --- $\Delta P = \frac{1}{2}\rho v^2$, $V = S\,\Delta P + V_0$ e
$ADC = V/q$ --- obtêm-se as constantes que descrevem apenas o sensor:

\begin{equation}
S = \frac{2\,q\,k}{\rho} = \SI{629,8}{\micro\volt\per\pascal}
\qquad
V_0 = \SI{2,66442}{\volt}
\label{eq:pitot_sensor}
\end{equation}

\noindent
Estas duas não mudam com a altitude, com a temperatura nem com o
microcontrolador, e são as efetivamente embarcadas na SIBA
(\autoref{sec:calibracao_pitot}). A \autoref{tab:pitot_conversores} mostra o
efeito de trocar o conversor: a SIBA lê com 12 bits e referência de
\SI{3,3}{\volt}, o que reduz o passo de quantização de \SI{4,89}{\milli\volt}
para \SI{0,806}{\milli\volt}.

\begin{table}[H]
    \centering
    \caption{Constantes equivalentes para diferentes conversores ($\rho = \SI{1,110}{\kilo\gram\per\meter\cubed}$)}
    \label{tab:pitot_conversores}
    \small
    \begin{tabular}{lcccc}
        \toprule
        \textbf{Configuração} & \textbf{$q$} & \textbf{$ADC_0$} &
        \textbf{$k$} & \textbf{1 LSB a \SI{5}{\meter\per\second}} \\
        & [\si{\micro\volt}/LSB] & [LSB] & [LSB/(\si{\meter\per\second})$^2$] &
        [\si{\meter\per\second}] \\
        \midrule
        10 bits / \SI{5,0}{\volt} & 4887,6 & 545,1  & 0,07152 & 1,398 \\
        10 bits / \SI{3,3}{\volt} & 3225,8 & 826,0  & 0,10837 & 0,923 \\
        12 bits / \SI{5,0}{\volt} & 1221,0 & 2182,2 & 0,28629 & 0,349 \\
        \textbf{12 bits / \SI{3,3}{\volt} (SIBA)} & \textbf{805,9} &
            \textbf{3306,3} & \textbf{0,43378} & \textbf{0,231} \\
        \bottomrule
    \end{tabular}

    \vspace{2pt}
    {\small Fonte: Autor (2026)}
\end{table}

\end{apendicesenv}
